\documentclass[12pt]{article}
\usepackage[margin=1.1in]{geometry}
\usepackage{titlesec}
\usepackage{xcolor}
\usepackage{hyperref}
\usepackage{enumitem}
\usepackage{fancyhdr}
\usepackage{setspace}
\usepackage{booktabs}
\usepackage[table]{xcolor}
\usepackage{fontspec}
\usepackage{enumitem}
\usepackage{amsmath}
\usepackage{newpxtext}
\usepackage{newpxmath}
\usepackage{comment}

%% colors in Duke's navy blue
\definecolor{titlecolor}{RGB}{0,0,128}

\hypersetup{
    colorlinks=true,
    linkcolor=titlecolor,
    urlcolor=titlecolor
}

%% itemize style
\setlist[itemize]{noitemsep, topsep=2pt}

%% section style
\titleformat{\section}{\color{titlecolor}\normalfont\Large\bfseries}{}{0em}{}
\titleformat{\subsection}{\normalfont\bfseries}{}{0em}{}
\setlength{\parskip}{3pt}

%% footer only (no top line)
\pagestyle{fancy}
\fancyhf{}
\renewcommand{\headrulewidth}{0pt}
\rfoot{\thepage}

%% subsection font size
\titleformat{\subsection}
  {\normalfont\fontsize{13.5}{15}\selectfont\bfseries}{}{0em}{}

\begin{document}

%% title block
\noindent {\Large \textcolor{titlecolor}{SOCIOL 723 Social Statistics II \\[3pt] Estimation, Causal Inference, and Machine Learning}}\\[0pt]

\subsection*{Time and Location}

Fall 2026\\
Tuesdays (Lecture) and Thursdays (Lab), 10:05AM--12:35PM\\
Erwin Square Plaza Room 753

\vspace{0.5em}

\subsection*{Instructor Information}

\begin{tabular}{@{}l l@{}}
\textbf{\textcolor{titlecolor}{Professor}} & Wenhao Jiang \\
\textbf{\textcolor{titlecolor}{Email}} & \href{mailto:wenhao.jiang@duke.edu}{wenhao.jiang@duke.edu} \\
\textbf{\textcolor{titlecolor}{Office Hours}} & Wednesdays 2--4PM or by appointment \\
\end{tabular}

\vspace{0.5em}

\subsection*{Course Overview and Prerequisites}

This is the second course in the graduate statistics sequence for sociology doctoral students. \textit{Social Statistics I}, using the model-comparison approach of Correll et al., rigorously developed regression, ANOVA, and a first look at logistic and Poisson regression in \texttt{R}. Building on that foundation, this course extends into more advanced quantitative methodology, organized around three pillars: estimation, causal inference, and machine learning. We begin by revisiting ordinary least squares (OLS), recasting the model-comparison framework in matrix form and attending carefully to its statistical properties, its interpretation as a tool for causal adjustment, and inference under realistic assumptions, including heteroskedasticity- and cluster-robust standard errors. We then develop maximum likelihood estimation (MLE) as the general engine behind most models sociologists use, grounding the logistic and Poisson models introduced in \textit{Social Statistics I} in likelihood theory and extending them, within the generalized linear model framework, to multinomial, ordered, and other limited dependent variables.

The core of the course turns to causal inference, the central methodological concern of contemporary empirical social science. We introduce the potential outcomes framework and directed acyclic graphs (DAGs) as complementary languages for defining and identifying causal effects, then work through the standard research designs: selection on observables (matching, propensity scores, and weighting), instrumental variables (IV), regression discontinuity designs (RDD), difference-in-differences (DID) with panel data, and synthetic control. The course closes by introducing machine learning for prediction and its integration with causal inference, including double/debiased machine learning and the estimation of heterogeneous treatment effects.

Each week pairs a Tuesday lecture that develops the theory with a Thursday lab that implements the methods on real data in \texttt{R}. By the end of the course, students should be able to read and critically evaluate the quantitative sociological literature that uses these methods, and to apply them competently in their own research. The course assumes working knowledge of the general linear model and statistical inference at the level of \textit{Social Statistics I}. Familiarity with matrix algebra and calculus is helpful but not required; the necessary tools are reviewed in the first two labs. \texttt{R} is the primary software; \texttt{Python} is introduced as a supplement for machine learning.

\subsection*{Expectations}

\textbf{\textcolor{titlecolor}{Reading and participation (10\%)}}:\\
You are expected to complete the required readings before the class in which they are discussed and to participate actively in both lecture and lab. If a reading is difficult, come prepared to discuss what you did and did not follow.

\vspace{0.5em}

\noindent \textbf{\textcolor{titlecolor}{Problem sets (30\%)}}:\\
There will be 5 problem sets, assigned roughly bi-weekly. They reinforce the lecture material and typically require \texttt{R} coding building on that week's lab. Problem sets are released on Tuesdays after class and are generally due two weeks later by Monday, 11:59PM (see the schedule for exact dates).

\vspace{0.5em}

\noindent \textbf{\textcolor{titlecolor}{In-Class Midterm Exam (25\%)}}: \\
There will be an in-class midterm on \textbf{Thursday, October 15}, covering the estimation half of the course (OLS and MLE) and the foundations of causal inference through instrumental variables (Weeks 1--7). The exam is open book and open note but timed. You will not be required to write any code.

\vspace{0.5em}

\noindent \textbf{\textcolor{titlecolor}{Final Project: Replication and Extension (35\%)}}: \\
In lieu of a final exam, you will complete an original project that replicates and extends a published study using one or more of the methods from the course, ideally connected to your own research. You should submit a short report (roughly 2,000--3,000 words) together with reproducible \texttt{R} code. Please sign up for the article you intend to replicate and confirm that replication data are available by \textbf{November 2}. The completed paper is due \textbf{December 14}. Using your own original data is also acceptable with prior approval.

\subsection*{Software}
This class uses \texttt{R}, and prior experience at the level of \textit{Social Statistics I} is assumed. \texttt{Python} will be used as a supplementary tool in certain cases, such as for machine learning. Please install a recent version of \texttt{R} and \texttt{RStudio} before the first lab, and use \texttt{RMarkdown} or \texttt{Quarto} for your assignments to support reproducibility.

\subsection*{Use of Artificial Intelligence (AI)}
While the use of AI is not permitted during the midterm exam, I encourage students to use AI tools to support their statistical learning, particularly for understanding steps in derivations and for debugging \texttt{R} code. AI assistance is also permitted on problem sets and the final project, but students are expected to fully understand every step of what they submit and to write up results in their own words.

\subsection*{Course Materials}

We will draw regularly on two required textbooks:
\begin{itemize}[noitemsep, topsep=2pt]
    \item \textcolor{titlecolor}{[MHE]} Angrist, Joshua D. and Jörn-Steffen Pischke. 2009. \textit{Mostly Harmless Econometrics: An Empiricist's Companion}. Princeton University Press. (our main reference for causal research designs)
    \item \textcolor{titlecolor}{[CCI]} Morgan, Stephen L. and Christopher Winship. 2015. \textit{Counterfactuals and Causal Inference: Methods and Principles for Social Research}, 2nd ed. Cambridge University Press. (the potential-outcomes and DAG framework)
\end{itemize}

The full text of \href{https://www.dsecoaching.com/pdf/2008%20Angrist%20Pischke%20MostlyHarmlessEconometrics.pdf}{\textcolor{titlecolor}{MHE}} is freely available online. The remaining references are used only for specific weeks and are freely available or optional rather than required purchases:

\begin{itemize}[noitemsep]
    \item \textcolor{titlecolor}{[ISL]} James, Gareth et al. 2021. \textit{An Introduction to Statistical Learning: with Applications in R}, 2nd ed. Springer. (\href{https://www.statlearning.com/}{online}; logistic and Poisson regression in Chapter 4, and the two machine-learning weeks)
    \item \textcolor{titlecolor}{[Pawitan]} Pawitan, Yudi. 2013. \textit{In All Likelihood: Statistical Modeling and Inference Using Likelihood}. Oxford University Press. (maximum likelihood theory)
    \item \textcolor{titlecolor}{[Mixtape]} Cunningham, Scott. 2021. \textit{Causal Inference: The Mixtape}. Yale University Press. (\href{https://mixtape.scunning.com/}{online})
    \item \textcolor{titlecolor}{[CML]} Chernozhukov, Victor et al. 2025. \textit{Applied Causal Inference Powered by ML and AI}. (\href{https://causalml-book.org/}{online}; machine-learning weeks)
\end{itemize}

For review of prerequisite material, students may consult the text used in \textit{Social Statistics I}: Correll, Josh, Abigail M. Folberg, Charles M. Judd, Gary H. McClelland, and Carey S. Ryan, \textit{Data Analysis: A Model Comparison Approach to Regression, ANOVA, and Beyond} (Routledge), together with its \texttt{R} companion, Stephen Vaisey's \href{https://vaiseys.github.io/damca/}{\textit{DAMCA}}.

\subsection*{Schedule}

The schedule is tentative as of \today. A detailed list of topics and readings, by week, appears below. Tuesday sessions are lectures and Thursday sessions are \texttt{R} labs. You are responsible for completing the \textit{required readings} before the Tuesday lecture.

\vspace{0.5em}

\begin{table}[h!]
\centering
\small
\setlength{\tabcolsep}{4pt}
\begin{tabular}{@{}llp{8.9cm}cc@{}}
\hline \hline\\[-10pt]
\textbf{\textcolor{black}{Week}} & \textbf{\textcolor{black}{Dates}} & \textbf{\textcolor{black}{Topic}} & \multicolumn{2}{c@{}}{\textbf{\textcolor{black}{Problem sets}}} \\
 & & & Assign & Due \\
\hline \\[-10pt]
1  & Aug 25/27  & The Linear Model and OLS &  &  \\
2  & Sep 1/3    & Regression Inference and Robust Standard Errors & 1 &  \\
3  & Sep 8/10   & Maximum Likelihood: Theory &  &  \\
4  & Sep 15/17  & Maximum Likelihood: Applications; Event History & 2 & 1 \\[4pt]
\arrayrulecolor{gray}
\hline \\[-9pt]
5  & Sep 22/24  & Causal Inference: Potential Outcomes and DAGs &  &  \\
6  & Sep 29/Oct 1 & Matching, Propensity Scores, and Weighting & 3 & 2 \\
7  & Oct 6/8    & Instrumental Variables &  &  \\
8  & Oct 13/15  & \textit{Fall break} (Tues) and \textit{In-class Midterm} (Thurs) &  &  \\
9  & Oct 20/22  & Regression Discontinuity Designs &  & 3 \\
10 & Oct 27/29  & Panel Data and Difference-in-Differences & 4  &  \\
11 & Nov 3/5    & Synthetic Control and Its Extensions &  &  \\[4pt]
\arrayrulecolor{gray}
\hline \\[-9pt]
12 & Nov 10/12  & Machine Learning: Prediction and Regularization & 5 & 4 \\
13 & Nov 17/19  & Causal Machine Learning & &  \\
14 & Nov 24     & Synthesis and Exam Review &  & 5 \\
   & Dec 3 & \textit{Cumulative final exam} (Thursday, 9:05--11:35AM) & &  \\
\arrayrulecolor{black}
\hline \hline
\end{tabular}
\flushleft {\footnotesize Tuesday sessions are lectures; Thursday sessions are hands-on \texttt{R} labs applying the week's methods. Fall break is Oct 9--14 (no Tuesday class Oct 13); the Thanksgiving recess begins Nov 24, so there is no Thursday lab on Nov 26.}
\end{table}

\noindent \textbf{\textcolor{titlecolor}{Week 1 (Aug 25/27): The Linear Model and OLS}}

\noindent \textit{Lecture}: The general linear model in scalar and matrix form, connecting the model-comparison framework of \textit{Social Statistics I} to matrix notation; the OLS estimator and its geometry; the Frisch--Waugh--Lovell theorem, omitted-variable bias, and regression as a tool for adjustment. \textit{Lab}: \texttt{R} and \texttt{RStudio} refresher; a review of the matrix algebra used in the linear model; fitting and interpreting OLS.
\begin{itemize}[noitemsep, topsep=2pt]
    \item[] \textit{Required}: MHE, Chapter 3; CCI, Chapter 6.
    \item[] \textit{Additional}: ISL, Chapters 2--3.
\end{itemize}

\vspace{0.8em} \noindent \textbf{\textcolor{titlecolor}{Week 2 (Sep 1/3): Review: Regression Inference and Robust Standard Errors}}

\noindent \textit{Lecture}: Sampling distribution of OLS and hypothesis testing; heteroskedasticity; robust and cluster-robust standard errors; the bootstrap; what robust standard errors can and cannot fix. \textit{Lab}: a review of the calculus used in estimation and optimization; robust and clustered inference and the bootstrap in \texttt{R}.
\begin{itemize}[noitemsep, topsep=2pt]
    \item[] \textit{Required}: MHE, Chapter 8.
    \item[] \textit{Additional}: Freedman, David A. 2006. ``On the So-Called `Huber Sandwich Estimator' and `Robust Standard Errors'.'' \textit{The American Statistician} 60(4), 299--302; Bertrand, Marianne, Esther Duflo, and Sendhil Mullainathan. 2004. ``How Much Should We Trust Differences-in-Differences Estimates?'' \textit{Quarterly Journal of Economics} 119(1), 249--275.
\end{itemize}

\vspace{0.8em} \noindent \textbf{\textcolor{titlecolor}{Week 3 (Sep 8/10): Maximum Likelihood: Theory}}

\noindent \textit{Lecture}: The likelihood and log-likelihood; the score function, Hessian, and Fisher information; identification and numerical maximization; the asymptotic properties of the MLE (consistency, asymptotic normality, and efficiency) under regularity conditions; and likelihood-ratio, Wald, and score tests. \textit{Lab}: coding a likelihood and optimizing it numerically in \texttt{R}.
\begin{itemize}[noitemsep, topsep=2pt]
    \item[] \textit{Required}: Pawitan, Chapters 2 and 4.
    \item[] \textit{Additional}: Pawitan, Chapter 3.
\end{itemize}

\vspace{0.8em} \noindent \textbf{\textcolor{titlecolor}{Week 4 (Sep 15/17): Maximum Likelihood: Applications}}

\noindent \textit{Lecture}: The generalized linear model framework (linear predictor, link, and likelihood); binary outcomes with logit and probit, now grounded in likelihood theory; interpretation via marginal effects, predicted probabilities, and interaction terms; multinomial and ordered models; count models (Poisson and negative binomial). \textit{Lab}: estimating and interpreting limited-dependent-variable models in \texttt{R}.
\begin{itemize}[noitemsep, topsep=2pt]
    \item[] \textit{Required}: ISL, Chapter 4; Hanmer, Michael J. and Kerem Ozan Kalkan. 2013. ``Behind the Curve: Clarifying the Best Approach to Calculating Predicted Probabilities and Marginal Effects from Limited Dependent Variable Models.'' \textit{American Journal of Political Science} 57(1), 263--277.
    \item[] \textit{Additional}: Berry, William D., Jacqueline H.R. DeMeritt, and Justin Esarey. 2010. ``Testing for Interaction in Binary Logit and Probit Models: Is a Product Term Essential?'' \textit{American Journal of Political Science} 54(1), 248--266; Zorn, Christopher. 2005. ``A Solution to Separation in Binary Response Models.'' \textit{Political Analysis} 13, 157--170.
\end{itemize}

\vspace{0.8em} \noindent \textbf{\textcolor{titlecolor}{Week 5 (Sep 22/24): Causal Inference: Potential Outcomes and DAGs}}

\noindent \textit{Lecture}: The potential outcomes framework; SUTVA, ignorability, and identification; average treatment effects; directed acyclic graphs, back-door and front-door criteria, and collider bias; defining a well-posed estimand. \textit{Lab}: simulating confounding and selection; drawing and reading DAGs in \texttt{R} (\texttt{dagitty}).
\begin{itemize}[noitemsep, topsep=2pt]
    \item[] \textit{Required}: MHE, Chapter 2; CCI, Chapters 1--3; Lundberg, Ian, Rebecca Johnson, and Brandon M. Stewart. 2021. ``What Is Your Estimand? Defining the Target Quantity Connects Statistical Evidence to Theory.'' \textit{American Sociological Review} 86(3), 532--565.
    \item[] \textit{Additional}: Holland, Paul W. 1986. ``Statistics and Causal Inference.'' \textit{Journal of the American Statistical Association} 81(396), 945--960; Greenland, Sander and Judea Pearl. 2017. ``Causal Diagrams.'' \textit{Wiley StatsRef: Statistics Reference Online}, 1--10.
\end{itemize}

\vspace{0.8em} \noindent \textbf{\textcolor{titlecolor}{Week 6 (Sep 29/Oct 1): Matching, Propensity Scores, and Weighting}}

\noindent \textit{Lecture}: Selection on observables; exact, nearest-neighbor, and coarsened matching; the propensity score; inverse-probability weighting; doubly robust estimation; ``bad controls'' and post-treatment bias. \textit{Lab}: matching and weighting workflows in \texttt{R} (\texttt{MatchIt}, \texttt{WeightIt}).
\begin{itemize}[noitemsep, topsep=2pt]
    \item[] \textit{Required}: CCI, Chapters 4, 5, 7, and 8.
    \item[] \textit{Additional}: Rosenbaum, Paul R. and Donald B. Rubin. 1983. ``The Central Role of the Propensity Score in Observational Studies for Causal Effects.'' \textit{Biometrika} 70(1), 41--55; Sekhon, Jasjeet S. 2009. ``Opiates for the Matches: Matching Methods for Causal Inference.'' \textit{Annual Review of Political Science} 12, 487--508.
\end{itemize}

\vspace{0.8em} \noindent \textbf{\textcolor{titlecolor}{Week 7 (Oct 6/8): Instrumental Variables}}

\noindent \textit{Lecture}: Endogeneity and the IV estimator; two-stage least squares; the LATE framework, monotonicity, and compliers; weak instruments and sensitivity. \textit{Lab}: 2SLS and diagnostics in \texttt{R} (\texttt{fixest}/\texttt{ivreg}).
\begin{itemize}[noitemsep, topsep=2pt]
    \item[] \textit{Required}: MHE, Chapter 4; CCI, Chapter 9.
    \item[] \textit{Additional}: Angrist, Joshua D., Guido W. Imbens, and Donald B. Rubin. 1996. ``Identification of Causal Effects Using Instrumental Variables.'' \textit{Journal of the American Statistical Association} 91(434), 444--455; Sovey, Allison J. and Donald P. Green. 2011. ``Instrumental Variables Estimation in Political Science: A Reader's Guide.'' \textit{American Journal of Political Science} 55(1), 188--200.
\end{itemize}

\vspace{0.8em} \noindent \textbf{\textcolor{titlecolor}{Week 8 (Oct 13/15): Fall Break / In-Class Midterm}}

\noindent \textit{Tuesday, Oct 13}: no class (Fall break, Oct 9--14). \textit{Thursday, Oct 15}: in-class midterm covering Weeks 1--7.

\vspace{0.8em} \noindent \textbf{\textcolor{titlecolor}{Week 9 (Oct 20/22): Regression Discontinuity Designs}}

\noindent \textit{Lecture}: Sharp and fuzzy RD; identification at the cutoff; local linear estimation, bandwidth choice, and kernel weighting; validity, sorting, and placebo checks. \textit{Lab}: RD estimation and diagnostics in \texttt{R} (\texttt{rdrobust}).
\begin{itemize}[noitemsep, topsep=2pt]
    \item[] \textit{Required}: MHE, Chapter 6; Imbens, Guido W. and Thomas Lemieux. 2008. ``Regression Discontinuity Designs: A Guide to Practice.'' \textit{Journal of Econometrics} 142(2), 615--635.
    \item[] \textit{Additional}: Lee, David S. and Thomas Lemieux. 2010. ``Regression Discontinuity Designs in Economics.'' \textit{Journal of Economic Literature} 48(2), 281--355; McCrary, Justin. 2008. ``Manipulation of the Running Variable in the Regression Discontinuity Design: A Density Test.'' \textit{Journal of Econometrics} 142(2), 698--714.
\end{itemize}

\vspace{0.8em} \noindent \textbf{\textcolor{titlecolor}{Week 10 (Oct 27/29): Panel Data and Difference-in-Differences}}

\noindent \textit{Lecture}: Panel data and unobserved heterogeneity; pooled OLS, first-difference, and fixed- and random-effects (within) estimators; two-way (unit and time) fixed effects (TWFE); the canonical 2$\times$2 DID and the parallel-trends assumption; DID and event studies as special cases of TWFE; problems with TWFE under staggered adoption and recent robust estimators. \textit{Lab}: panel and DID estimation in \texttt{R} (\texttt{plm}, \texttt{fixest}, \texttt{did}).
\begin{itemize}[noitemsep, topsep=2pt]
    \item[] \textit{Required}: MHE, Chapter 5; CCI, Chapter 11; Mixtape, Difference-in-Differences chapter.
    \item[] \textit{Additional}: Goodman-Bacon, Andrew. 2021. ``Difference-in-Differences with Variation in Treatment Timing.'' \textit{Journal of Econometrics} 225(2), 254--277; Roth, Jonathan et al. 2023. ``What's Trending in Difference-in-Differences? A Synthesis of the Recent Econometrics Literature.'' \textit{Journal of Econometrics} 235(2), 2218--2244.
\end{itemize}

\vspace{0.8em} \noindent \textbf{\textcolor{titlecolor}{Week 11 (Nov 3/5): Synthetic Control and Event Studies}}

\noindent \textit{Lecture}: Comparative case studies; the synthetic control estimator; inference with few treated units; generalized synthetic control and event-study designs. \textit{Lab}: synthetic control in \texttt{R} (\texttt{Synth}, \texttt{gsynth}).
\begin{itemize}[noitemsep, topsep=2pt]
    \item[] \textit{Required}: Abadie, Alberto, Alexis Diamond, and Jens Hainmueller. 2010. ``Synthetic Control Methods for Comparative Case Studies: Estimating the Effect of California's Tobacco Control Program.'' \textit{Journal of the American Statistical Association} 105(490), 493--505.
    \item[] \textit{Additional}: Abadie, Alberto and Javier Gardeazabal. 2003. ``The Economic Costs of Conflict: A Case Study of the Basque Country.'' \textit{American Economic Review} 93(1), 113--132; Xu, Yiqing. 2017. ``Generalized Synthetic Control Method: Causal Inference with Interactive Fixed Effects Models.'' \textit{Political Analysis} 25(1), 57--76.
\end{itemize}

\vspace{0.8em} \noindent \textbf{\textcolor{titlecolor}{Week 12 (Nov 10/12): Machine Learning: Prediction and Regularization}}

\noindent \textit{Lecture}: Prediction versus inference; the bias--variance tradeoff; cross-validation; ridge, lasso, and elastic net; trees, random forests, and boosting. \textit{Lab}: regularized regression and tree ensembles in \texttt{R} (\texttt{glmnet}, \texttt{ranger}).
\begin{itemize}[noitemsep, topsep=2pt]
    \item[] \textit{Required}: ISL, Chapters 5, 6, and 8; Molina, Mario and Filiz Garip. 2019. ``Machine Learning for Sociology.'' \textit{Annual Review of Sociology} 45(1), 27--45.
    \item[] \textit{Additional}: ISL, Chapter 10 (deep learning); Kleinberg, Jon et al. 2015. ``Prediction Policy Problems.'' \textit{American Economic Review: Papers \& Proceedings} 105(5), 491--495.
\end{itemize}

\vspace{0.8em} \noindent \textbf{\textcolor{titlecolor}{Week 13 (Nov 17/19): Causal Machine Learning}}

\noindent \textit{Lecture}: Why naive ML is biased for causal parameters; Neyman orthogonality and double/debiased machine learning (DML); partialling-out and cross-fitting; heterogeneous treatment effects and causal forests. \textit{Lab}: DML and heterogeneous effects in \texttt{R} (\texttt{DoubleML}, \texttt{grf}).
\begin{itemize}[noitemsep, topsep=2pt]
    \item[] \textit{Required}: CML, Chapters 4 and 10; Athey, Susan and Guido W. Imbens. 2017. ``The State of Applied Econometrics: Causality and Policy Evaluation.'' \textit{Journal of Economic Perspectives} 31(2), 3--32.
    \item[] \textit{Additional}: Belloni, Alexandre, Victor Chernozhukov, and Christian Hansen. 2014. ``High-Dimensional Methods and Inference on Structural and Treatment Effects.'' \textit{Journal of Economic Perspectives} 28(2), 29--50; Wager, Stefan and Susan Athey. 2018. ``Estimation and Inference of Heterogeneous Treatment Effects Using Random Forests.'' \textit{Journal of the American Statistical Association} 113(523), 1228--1242.
\end{itemize}

\vspace{0.8em} \noindent \textbf{\textcolor{titlecolor}{Week 14 (Nov 24): Synthesis and Final Project Workshop}}

\noindent \textit{Tuesday, Nov 24}: a synthesis of how estimation, causal identification, and machine learning fit together, followed by an in-class workshop of final projects. \textit{Thursday, Nov 26}: no class (Thanksgiving recess). The final project is due \textbf{December 14}.

\end{document}